% Exact excerpts extracted from the frozen Lean sources; no proof body is reproduced.
% Preamble requirements: \usepackage{fontspec,fancyvrb}
%                       \newfontfamily\leanappendixfont{DejaVuSansMono.ttf}
% Insert after \appendix in the main document. This snippet adds no page breaks.
% Source SHA256 ReversedBinaryPowerClosure.lean: 6f448bb1475e4214a10d4acdac8835d526a559d166c93acfd910e39bca23f281
% Source SHA256 ReversedRealNormalization.lean: 947bb86b18f46ba0f150437c4a334c79a18b89bbc73b5855eff9381237b23761
% Source SHA256 FullTwoBasic.lean: 38d6e8029e7fff5a948aca33ad661b28ed232ac527c3627882eac397cf00780e
% Source SHA256 CollatzReversedRealTwoCoordinate.lean: b7cc9035e073d624fd270f6fbde9f2149a290dd93f238643a65aefe4fe434b71
% Source SHA256 FullTwoCoordinate.lean: 429785555eecec277248380fe0c2e767c8c7303cf0e2712de29878dda2bd29a0

\section{Exact Lean declarations}
\label{app:lean-statements}
The blocks below reproduce the checked source verbatim, with the original
line numbers. The theorem blocks end at their types; their proof terms are
omitted. Source paths are relative to the repository root.

The affine definitions and \texttt{ReversedRealWeak} are in
\texttt{CollatzCertificate}. The definitions and theorem with diagonal bounds
are in \texttt{CollatzResearch.FullTwo}. Matrix notation is opened in both
namespaces. The order on vectors is pointwise, so \texttt{Weak X Y} means
that every coefficient of \texttt{Y} is at most the corresponding coefficient
of \texttt{X}. Composition places the left argument outside the right argument.
The unrestricted reversed theorem is in
\texttt{CollatzResearch.RealTwoCoordinate}; it uses only
\texttt{Affine.Nonnegative} and the eleven reversed weak comparisons.

\paragraph{Ambient types and the coefficient order.}
\noindent\texttt{\detokenize{formal/ReversedBinaryPowerClosure.lean}}, lines 10--17.
\begin{Verbatim}[formatcom=\leanappendixfont,fontsize=\scriptsize,numbers=left,firstnumber=10,numbersep=5pt,xleftmargin=18pt]
variable {n : Type*} [Fintype n]

abbrev Vec (n : Type*) := n → ℝ
abbrev Mat (n : Type*) := Matrix n n ℝ

def EntrywiseLE (M N : Mat n) : Prop := ∀ i j, M i j ≤ N i j

local infix:50 " ≤ₘ " => EntrywiseLE
\end{Verbatim}

\paragraph{Affine composition and comparison.}
\noindent\texttt{\detokenize{formal/ReversedBinaryPowerClosure.lean}}, lines 255--266.
\begin{Verbatim}[formatcom=\leanappendixfont,fontsize=\scriptsize,numbers=left,firstnumber=255,numbersep=5pt,xleftmargin=18pt]
structure Affine (n : Type*) where
  matrix : Mat n
  offset : Vec n

def Affine.comp (X Y : Affine n) : Affine n :=
  ⟨X.matrix * Y.matrix, X.matrix *ᵥ Y.offset + X.offset⟩

def Affine.Nonnegative (X : Affine n) : Prop :=
  (0 ≤ₘ X.matrix) ∧ 0 ≤ X.offset

def Affine.Weak (X Y : Affine n) : Prop :=
  (Y.matrix ≤ₘ X.matrix) ∧ Y.offset ≤ X.offset
\end{Verbatim}

\paragraph{Two coordinates and admissibility.}
\noindent\texttt{\detokenize{formal/FullTwoBasic.lean}}, lines 9--11.
\begin{Verbatim}[formatcom=\leanappendixfont,fontsize=\scriptsize,numbers=left,firstnumber=9,numbersep=5pt,xleftmargin=18pt]
abbrev Aff2 := Affine (Fin 2)

def Admissible (X : Aff2) : Prop := X.Nonnegative ∧ 1 ≤ X.matrix 0 0
\end{Verbatim}

\paragraph{Forward weak rules.}
\noindent\texttt{\detokenize{formal/FullTwoBasic.lean}}, lines 21--32.
\begin{Verbatim}[formatcom=\leanappendixfont,fontsize=\scriptsize,numbers=left,firstnumber=21,numbersep=5pt,xleftmargin=18pt]
structure ForwardWeak (A B C D E F G : Aff2) : Prop where
  ad : (A.comp D).Weak D
  bd : (B.comp D).Weak (G.comp D)
  ae : (A.comp E).Weak (E.comp A)
  af : (A.comp F).Weak (E.comp B)
  ag : (A.comp G).Weak (F.comp A)
  be : (B.comp E).Weak (F.comp B)
  bf : (B.comp F).Weak (G.comp A)
  bg : (B.comp G).Weak (G.comp B)
  ce : (C.comp E).Weak (C.comp B)
  cf : (C.comp F).Weak (C.comp (A.comp A))
  cg : (C.comp G).Weak (C.comp (A.comp B))
\end{Verbatim}

\paragraph{Forward offset equalities.}
\noindent\texttt{\detokenize{formal/FullTwoBasic.lean}}, lines 34--45.
\begin{Verbatim}[formatcom=\leanappendixfont,fontsize=\scriptsize,numbers=left,firstnumber=34,numbersep=5pt,xleftmargin=18pt]
structure ForwardGapsZero (A B C D E F G : Aff2) : Prop where
  ad : (A.comp D).offset 0 = D.offset 0
  bd : (B.comp D).offset 0 = (G.comp D).offset 0
  ae : (A.comp E).offset 0 = (E.comp A).offset 0
  af : (A.comp F).offset 0 = (E.comp B).offset 0
  ag : (A.comp G).offset 0 = (F.comp A).offset 0
  be : (B.comp E).offset 0 = (F.comp B).offset 0
  bf : (B.comp F).offset 0 = (G.comp A).offset 0
  bg : (B.comp G).offset 0 = (G.comp B).offset 0
  ce : (C.comp E).offset 0 = (C.comp B).offset 0
  cf : (C.comp F).offset 0 = (C.comp (A.comp A)).offset 0
  cg : (C.comp G).offset 0 = (C.comp (A.comp B)).offset 0
\end{Verbatim}

\paragraph{Reversed weak rules.}
Here the index type is finite and has decidable equality; the final
theorem instantiates it with \texttt{Fin 2}.
\noindent\texttt{\detokenize{formal/ReversedRealNormalization.lean}}, lines 83--94.
\begin{Verbatim}[formatcom=\leanappendixfont,fontsize=\scriptsize,numbers=left,firstnumber=83,numbersep=5pt,xleftmargin=18pt]
structure ReversedRealWeak (A B C D E F G : Affine ι) : Prop where
  da : (D.comp A).Weak D
  db : (D.comp B).Weak (D.comp G)
  ea : (E.comp A).Weak (A.comp E)
  fa : (F.comp A).Weak (B.comp E)
  ga : (G.comp A).Weak (A.comp F)
  eb : (E.comp B).Weak (B.comp F)
  fb : (F.comp B).Weak (A.comp G)
  gb : (G.comp B).Weak (B.comp G)
  ec : (E.comp C).Weak (B.comp C)
  fc : (F.comp C).Weak (A.comp (A.comp C))
  gc : (G.comp C).Weak (B.comp (A.comp C))
\end{Verbatim}

\paragraph{Reversed offset equalities.}
\noindent\texttt{\detokenize{formal/FullTwoBasic.lean}}, lines 47--58.
\begin{Verbatim}[formatcom=\leanappendixfont,fontsize=\scriptsize,numbers=left,firstnumber=47,numbersep=5pt,xleftmargin=18pt]
structure ReversedGapsZero (A B C D E F G : Aff2) : Prop where
  da : (D.comp A).offset 0 = D.offset 0
  db : (D.comp B).offset 0 = (D.comp G).offset 0
  ea : (E.comp A).offset 0 = (A.comp E).offset 0
  fa : (F.comp A).offset 0 = (B.comp E).offset 0
  ga : (G.comp A).offset 0 = (A.comp F).offset 0
  eb : (E.comp B).offset 0 = (B.comp F).offset 0
  fb : (F.comp B).offset 0 = (A.comp G).offset 0
  gb : (G.comp B).offset 0 = (B.comp G).offset 0
  ec : (E.comp C).offset 0 = (B.comp C).offset 0
  fc : (F.comp C).offset 0 = (A.comp (A.comp C)).offset 0
  gc : (G.comp C).offset 0 = (B.comp (A.comp C)).offset 0
\end{Verbatim}

\paragraph{The theorem with diagonal bounds.}
\noindent\texttt{\detokenize{formal/FullTwoCoordinate.lean}}, lines 30--34.
\begin{Verbatim}[formatcom=\leanappendixfont,fontsize=\scriptsize,numbers=left,firstnumber=30,numbersep=5pt,xleftmargin=18pt]
theorem full_two_coordinate_obstruction (A B C D E F G : Aff2)
    (hA : Admissible A) (hB : Admissible B) (hC : Admissible C) (hD : Admissible D)
    (hE : Admissible E) (hF : Admissible F) (hG : Admissible G) :
    (ForwardWeak A B C D E F G → ForwardGapsZero A B C D E F G) ∧
    (ReversedRealWeak A B C D E F G → ReversedGapsZero A B C D E F G)
\end{Verbatim}

\paragraph{The unrestricted reversed theorem.}
\noindent\texttt{\detokenize{formal/CollatzReversedRealTwoCoordinate.lean}}, lines 30--35.
\begin{Verbatim}[formatcom=\leanappendixfont,fontsize=\scriptsize,numbers=left,firstnumber=30,numbersep=5pt,xleftmargin=18pt]
theorem reversed_eligible_offsets_equal
    (A B C D E F G : Affine (Fin 2))
    (hA : A.Nonnegative) (hB : B.Nonnegative) (hC : C.Nonnegative)
    (hD : D.Nonnegative) (hE : E.Nonnegative) (hF : F.Nonnegative) (hG : G.Nonnegative)
    (h : ReversedRealWeak A B C D E F G) :
    (D.comp A).offset = D.offset ∧ (D.comp B).offset = (D.comp G).offset
\end{Verbatim}
